\chapter{Lipschitz extension and almost-everywhere differentiation}\label{sec:lipschitz-tools}
This section develops the analytic input for rectifiable parametrizations. The source is \cite{Sim18}, Chapter 2, Section 1, pp.~43--46. The one-variable fundamental theorem for absolutely continuous functions is a background input: such functions have an integrable derivative almost everywhere and are recovered by integrating it. Fubini and dominated convergence are also used explicitly. The many-variable differentiability argument is supplied below.

\begin{definition}[Lipschitz constants and local Lipschitz regularity]\label{def:lipschitz-constants}
A map $f:X\to Y$ between metric spaces is $L$-Lipschitz if $d_Y(f(x),f(y))\leq Ld_X(x,y)$ for all $x,y\in X$. Its least such constant is $\operatorname{Lip}f$. On a Euclidean open set, locally Lipschitz means every point has a neighborhood on which some finite constant works. The local constants need not be uniformly bounded on the whole open set. See \cite{Sim18}, equation (1.1), p.~43.
\end{definition}

\begin{proposition}[Scalar Lipschitz extension by cones]\label{prop:mcshane-extension}
Let $A$ be a nonempty subset of a metric space $X$, and let $f:A\to\mathbb R$ be $L$-Lipschitz. Then
\[
 F(x)=\inf_{a\in A}\{f(a)+Ld(x,a)\}
\]
is finite, is $L$-Lipschitz on $X$, and agrees with $f$ on $A$. If $M=\sup_A|f|<\infty$, truncating $F$ to $[-M,M]$ gives an extension with the same Lipschitz bound and the same supremum norm as $f$. This is \cite{Sim18}, Theorem 1.2, p.~44.
\end{proposition}
\begin{proof}
Fix $a_0\in A$. The triangle inequality and the Lipschitz bound give
\[
 f(a_0)-Ld(x,a_0)\leq f(a)+Ld(x,a)
\]
for every $a\in A$, while the candidate $a=a_0$ gives a finite upper bound for the infimum. If $x\in A$, every candidate is at least $f(x)$ and $a=x$ gives equality. For any $x,y$, each candidate satisfies $f(a)+Ld(x,a)\leq f(a)+Ld(y,a)+Ld(x,y)$. Take infima and exchange $x,y$ to prove the Lipschitz bound. Scalar truncation is 1-Lipschitz and fixes all values on $A$. Agreement on $A$ gives equality, rather than just an upper bound, of the supremum norms. No vector-valued extension with the identical constant is asserted by this scalar formula.
\end{proof}

\begin{lemma}[Lipschitz functions on lines are absolutely continuous]\label{lem:lipschitz-on-lines}
An $L$-Lipschitz real function on an interval is absolutely continuous on every compact subinterval, is differentiable almost everywhere, and its derivative has absolute value at most $L$. A locally Lipschitz real function on an open subset of $\mathbb R^m$ therefore has directional derivatives in any one fixed direction at almost every point locally.
\end{lemma}
\begin{proof}
For finitely many disjoint subintervals $(a_j,b_j)$, the sum of the absolute increments is at most $L\sum_j(b_j-a_j)$. This is the definition of absolute continuity. The one-variable fundamental theorem gives almost-everywhere differentiability, and each difference quotient is bounded by $L$. For the higher-dimensional assertion, restrict to a box lying in a neighborhood with a fixed Lipschitz constant, rotate the box to the chosen direction, and apply the one-dimensional assertion on its parallel lines. Fubini gives a null exceptional set in the box. Countably many such local boxes cover the domain. The nondifferentiability set is measurable since the limsup and liminf of difference quotients may be taken along rational increments. This is the line argument in \cite{Sim18}, proof of Theorem 1.4, p.~45.
\end{proof}

\begin{lemma}[Vanishing distribution tests determine a locally integrable function]\label{lem:zero-distribution-function}
If $h\in L^1_{\rm loc}(U)$ and $\int_Uh\phi=0$ for every $\phi\in C_c^\infty(U)$, then $h=0$ almost everywhere.
\end{lemma}
\begin{proof}
Fix a smooth nonnegative mollifier of integral one. At a point $x$ whose distance from the complement of $U$ exceeds $\varepsilon$, the function $y\mapsto\rho_\varepsilon(x-y)$ is an admissible test, so the local mollification is zero. At almost every point, Theorem \ref{thm:lebesgue-differentiation} and Proposition \ref{prop:mollifiers-at-lebesgue-points}, localized inside $U$, say that these mollifications tend to $h(x)$. Thus $h$ vanishes almost everywhere.
\end{proof}

\begin{lemma}[Directional derivatives are linear almost everywhere in each fixed direction]\label{lem:directional-weak-linearity}
Suppose $f:U\to\mathbb R$ is locally Lipschitz, and put $g=(\partial_1f,\ldots,\partial_mf)$ where the coordinate derivatives exist, giving them arbitrary values on the null exceptional set. Then $g\in L^\infty_{\rm loc}(U;\mathbb R^m)$. For each fixed unit vector $v$, the directional derivative $D_vf$ exists almost everywhere and equals $v\cdot g$ almost everywhere. The exceptional set in this assertion may depend on $v$.
\end{lemma}
\begin{proof}
Existence and local bounds follow from Lemma \ref{lem:lipschitz-on-lines}. If $\phi\in C_c^\infty(U)$, small difference quotients have a common compact support after multiplication by $\phi$ and are uniformly bounded there by local Lipschitz constants. Translation and dominated convergence give
\[
 \int D_vf\,\phi=-\int f\,D_v\phi.
\]
Use the same identity for the coordinate directions and $D_v\phi=\sum_iv_i\partial_i\phi$ to obtain $\int(D_vf-v\cdot g)\phi=0$. Lemma \ref{lem:zero-distribution-function} proves the claimed equality. Localizing the test to finitely many small neighborhoods justifies the bound when no single Lipschitz constant has been specified globally. Compare \cite{Sim18}, proof of Theorem 1.4, equations (2)--(4), p.~45.
\end{proof}

\begin{lemma}[A finite angular net upgrades directional limits]\label{lem:angular-net-differentiability}
Let $f$ be $L$-Lipschitz near $x\in\mathbb R^m$. Suppose $a\in\mathbb R^m$ and a dense set $D\subset S^{m-1}$ satisfy
\[
 \lim_{t\to0}\frac{f(x+tv)-f(x)}t=a\cdot v\qquad(v\in D).
\]
Then $f$ is differentiable at $x$ with derivative $h\mapsto a\cdot h$, and $|a|\leq L$.
\end{lemma}
\begin{proof}
The directional difference quotients are bounded by $L$. Taking limits on the dense set and then using continuity in $v$ gives $|a|\leq L$. Put $Q(v,t)=(f(x+tv)-f(x))/t-a\cdot v$. For fixed nonzero small $t$,
\[
 |Q(v,t)-Q(w,t)|\leq(L+|a|)|v-w|\leq2L|v-w|.
\]
Given $\eta>0$, compactness of the sphere supplies a finite $\eta$-net $v_1,\ldots,v_N$ from $D$. For all sufficiently small $|t|$, every $|Q(v_i,t)|<\eta$. Thus $\sup_v|Q(v,t)|\leq(2L+1)\eta$. Let $\eta\downarrow0$ and write an arbitrary nonzero increment as $h=|h|v$. The resulting uniform estimate is exactly differentiability. This is the finite-net step in \cite{Sim18}, p.~46.
\end{proof}

\begin{proposition}[Smooth approximation preserves a scalar Lipschitz bound]\label{prop:lipschitz-mollification}
Let $f:\mathbb R^m\to\mathbb R$ be $L$-Lipschitz, and let $\rho\in C_c^\infty(B_1)$ be nonnegative with integral one. Its mollifications satisfy
\[
 f_\varepsilon\in C^\infty,\qquad
 |f_\varepsilon-f|\leq L\varepsilon\int |z|\rho(z)\,dz,
 \qquad \operatorname{Lip}f_\varepsilon\leq L.
\]
Moreover $Df_\varepsilon\to Df$ almost everywhere, with the derivative interpreted by Theorem \ref{thm:rademacher}.
\end{proposition}
\begin{proof}
The kernel has compact support, so local integrability suffices for smoothness. Write $f_\varepsilon(x)=\int f(x-\varepsilon z)\rho(z)\,dz$ and apply the Lipschitz inequality, first to $x$ and $x-\varepsilon z$ and then to $x-\varepsilon z$ and $y-\varepsilon z$. This gives both bounds. The integration-by-parts identity in Lemma \ref{lem:directional-weak-linearity} shows that $\partial_i f_\varepsilon=\rho_\varepsilon*\partial_i f$. Lebesgue differentiation and convergence at Lebesgue points apply to these locally bounded weak derivatives. The coordinate limits agree with the derivative from Rademacher almost everywhere. The Lipschitz bound alone does not imply convergence of derivatives at every point.
\end{proof}

\begin{example}[All directional derivatives need not give differentiability]\label{ex:directional-not-differentiable}
Define $f(0,0)=0$ and $f(x,y)=x^3/(x^2+y^2)$ away from the origin. Then $f$ is globally Lipschitz and every directional derivative at zero exists, but $f$ is not differentiable there.
\end{example}
\begin{proof}
Away from zero,
\[
 f_x=\frac{x^4+3x^2y^2}{(x^2+y^2)^2},\qquad
 f_y=\frac{-2x^3y}{(x^2+y^2)^2}.
\]
These derivatives are bounded uniformly. Integrating the gradient on a segment gives a uniform Lipschitz bound if the segment avoids zero. If it passes through zero, split the segment and use continuity at zero, which follows from $|f(x,y)|\leq\sqrt{x^2+y^2}$. Homogeneity gives $D_vf(0)=v_1^3/(v_1^2+v_2^2)$ for every nonzero direction $v$. The coordinate derivatives are 1 and 0, so a derivative would have to be $(h_1,h_2)\mapsto h_1$. At $v=(1,1)$ the directional derivative is $1/2$, not 1. The directional values fail the linearity hypothesis of Lemma \ref{lem:angular-net-differentiability}.
\end{proof}
